\uuid{mz1r}
\titre{Somme/Produit/DL usuels}
\auteur{hanani ; mzari}
\organisation{exo7}
\licence{CC BY-NC-SA 4.0}
\datecreate{2020-01-20}
\niveau{L1}
\module{Analyse}
\chapitre{Développement limité}
\sousChapitre{Calculs}
\difficulte{1}
\theme{Opérations sur les DL, Somme, Produit, DL usuels}
\exo7qcmid{lille-2.developpements-limites.002}

\contenu{
\texte{
Soit $a\in \Rr$ et $f$ une fonction. On dit que $f(x)$ est un $o\left((x-a)^n\right)$ au voisinage de $a$ et on écrit $f(x)=o((x-a)^n)$ si $\displaystyle \lim _{x\to a}\frac{f(x)}{(x-a)^n}=0$. On note par $DL_n(a)f(x)$ le développement limité de $f$ à l'ordre $n$ en $a$. On note aussi $DL_n(a^+)f(x)$ (resp. $DL_n(a^-)f(x)$) le développement limité de $f$ à l'ordre $n$ à droite en $a$ (resp. à gauche en $a$).
}

\begin{question}
\qclassification{125.02}
\qtags{motcle=Somme/Produit/DL usuels}
Soient $f$ et $g$ deux fonctions telles que :
$$DL_3(0)f(x)=x+x^3+o(x^3)\; \mbox{ et }\; DL_3(0)g(x)=-x+x^3+o(x^3).$$
On peut en déduire que :
\begin{answers}
\good{$DL_2(0)\left[f(x)+g(x)\right]=o(x^2)$.}
\bad{$DL_2(0)\left[f(x)-g(x)\right]=o(x^2)$.}
\good{$DL_2(0)\left[2f(x)+g(x)\right]=x+o(x^2)$.}
\bad{$\displaystyle DL_6(0)f(x)\times g(x)=-x^2+x^6+o(x^6)$.}
\end{answers}
\begin{explanations}
Découle des opérations sur les DL en remarquant que la somme ou le produit de deux DL du même ordre donne un DL du même ordre.
\end{explanations}
\end{question}
}
